
\subsubsection*{Keep alternative annual histories and their statements together}
The same opening book can support several explicitly assumed sequences. A
company may have payment authority in one scenario and lack it in another;
the tax rate or recovery capacity may remain unknown. Comparing those
sequences requires keeping each company's balances, profit and cash movements
together. A cash balance from one scenario cannot be combined with the profit
or liabilities of another and presented as one company result.

An annual family names its members, states their assumptions and supplies one
common opening request and a consecutive reporting horizon. Each member lists
its own annual operations and planned period carries. The opening is bound by
its digest, and the planned periods must agree. Structural validation checks
these declarations without executing the financial methods. An unknown rate
can therefore remain in an otherwise well-formed member; the financial
evaluation determines where that member can proceed.

Each evaluated member uses the same annual evaluator as an individually
submitted program. Its complete financial history is retained. In the worked
sequence, the first period ends with cash $701/2$, recorded profit $-5$ and
operating cash flow $-7$. Its next period records profit $-2$, no cash movement
and total tax payable $27/4$. A fresh deferred measurement and another
provisional carry retain cash $701/2$ and change the debit-positive retained
balance to $81/4$. Both earlier periods remain available in the archive.

Other members need not reach those periods. If payment authority is false,
the first assessment succeeds but the payment is refused. The assessment,
its resulting balances and every remaining operation stay in that member's
result. Its later period has no report because execution never reached it.
An unknown rate stops before the first assessment; missing recovery capacity
or a stale measurement stops later. An explicit payment exceeding cash is
retained as a funding failure. None of these refusals removes the member or
prevents an independent member from continuing.

A joint period query returns the whole report for each member, bound to its
complete result. A period crossed by a successful provisional carry uses the
archived report. The current period uses its completed-prefix report and
states the actual date reached and the current and deferred assessment
statuses. Completing all listed operations does not turn that prefix into
an admitted annual statement. A period not reached is explicitly unavailable,
rather than a report of zeros. The query recomputes the family before trusting
saved results; changing cash and recomputing a content hash does not validate
an altered financial history.

The caller supplies an explicit member-evaluation budget and traversal order.
When that budget is exhausted, every unvisited member keeps its identity,
assumptions and request digest, with no invented financial result. Even a zero
member budget checks the common opening; that check is outside the count of
member evaluations. The count is not a guarantee about execution time or
memory. Increasing it makes a new request for recomputation, not a durable
continuation receipt.

These members are unweighted authored alternatives. Two names with different
stated assumptions remain distinct even when their financial results agree.
An empty authored selection does not establish that no feasible company
history exists. Evaluating every supplied member establishes coverage of that
list only. Reconstructing actual histories, proving an exhaustive feasible set
or attaching a probability law requires additional evidence and methods.
